\documentclass{article}
\usepackage[T1]{fontenc}
\usepackage{lmodern}
\usepackage{graphicx}
\usepackage{amsmath,amssymb}
\usepackage[a4paper,margin=2cm]{geometry}
\usepackage[absolute,overlay]{textpos}
\setlength{\TPHorizModule}{1mm}
\setlength{\TPVertModule}{1mm}
\usepackage[indonesian]{babel}
\usepackage{hyperref}
\setlength{\emergencystretch}{2em}
% unit: Halaman 2

\begin{document}

% === Halaman 2 (python/native) ===

2

Cipher Blok (Block Cipher)

• Berbeda dengan cipher alir, pada cipher blok plainteks dibagi menjadi blok-blok bit 

dengan panjang sama. Ukuran blok yang umum adalah 64 bit, 128 bit,  256 bit, dsb.

• Panjang blok cipherteks = panjang blok plainteks.

• Enkripsi dilakukan pada setiap blok plainteks dengan menggunakan bit-bit kunci 

• Panjang kunci eksternal (yang diberikan oleh pengguna) tidak harus sama dengan 

panjang blok plainteks. Pada beberapa cipher blok, panjang kunci = panjang blok, 

tetapi pada cipher blok yang lain tidak sama.

\end{document}
